\subsection{Outline of Questions}
\label{appendix-outline}
Here we provide an outline of the interview questions used. Within each interview topic (e.g. Introduction), we show examples of some of the different questions that were asked based on the nature of the conversation. Given the semi-structured nature, we did not adhere to the sequence of these questions within each topic. 

\subsubsection{Introduction}
Briefly introduce the study's focus on multi-sensory biophilic design within smart buildings highlighting the interest in MSB around.
\begin{enumerate}
    \item Could you share your background in (relevant text)---particularly regarding your experience with biophilic design? 
    \item Interpretation of Connection: When people talk about this goal for design that allows “a deeper connection to nature” in built environments, what do you think that means?
    \item How does this resonate with your experiences or design philosophy?
    \item (give context on smart environments) What are your thoughts on integrating biophilic elements within smart environments?
\end{enumerate}


\subsubsection{Grounded Insights} Explore the practical and technical aspects, identifying research-based perspectives, technical insights, criticisms, and relevant examples in MSB design.
\begin{enumerate}
    \item Existing Design Limitations: (...our desk research shows...) Current biophilic designs are often visually focused (e.g. use of lights, living walls, nature-inspired privacy screens, green colours). Do you agree, and why do you think that is the case? 
    \item Sensory Engagement: Which sensory elements (e.g. smell, texture, sounds, even vibrations) do you think could be most effective in creating a sense of immersion or connection to nature within built environments?
\end{enumerate}


\subsubsection{Challenges and Feasibility}
\begin{enumerate}
    \item What [technical] challenges do you anticipate when designing multi-sensory biophilic spaces, such as those incorporating sound, touch, or scent? (depending on the background - focus on either future material possibilities, future concepts, or technology)
    \item Speculating on Future Biophilic Technologies: Are there any emerging technologies or materials that you think could enhance multi-sensory biophilic design? How might these help create environments that are more ‘alive’ and responsive?
    \item Criticisms: Are there any problems you have with the way we consider nature design in its current applications?
\end{enumerate}


\subsubsection{Experiential Perspectives} 
Gather insights on personal experiences, professional values, and context-specific knowledge regarding MSB design.
\begin{enumerate}
    \item Body-Nature Connection: Can you share an experience or a project where a built environment made you feel deeply connected to nature? Which aspects were most impactful?
    \item Multi-sensory: Similarly, are there any specific projects you've worked on or encountered that successfully engaged multiple senses? Do you remember what impact it had on you? 
    \item Values and Priorities: What core values (e.g. peace, healing, diversity) do you think should shape the design of a MSB space?
\end{enumerate}


\subsubsection{Visionary Outlooks} Encourage exploration of  future visions and speculate on potential evolutions in MSB design.
\begin{enumerate}
    \item Exploring Enhancement Possibilities: “If you could enhance one sensory aspect of biophilic design in current spaces, what would it be, and how do you imagine it might affect occupants’ experiences?”
    \item  Speculative Scenarios: “Picture a future environment that seamlessly blends technology with natural elements—for instance, a space that feels like it ‘breathes’ (inhales and exhales). Are there any such forward-thinking or aspirational concepts you've come across—or envisioned yourself—that could inform the future of multi-sensory environments? 
    \item Future Visions: If you could design a biophilic environment that integrates various sensory responses to create a deep connection with nature, what would this look and feel like? 
\end{enumerate}

\subsubsection{Conclusion and Reflection}
Wrap up with open reflection and any additional thoughts.
\begin{enumerate}
    \item Do you see biophilic design gaining more importance in the future where spaces are increasingly smart? 
    \item What role do you see for architects/designers in shaping these sensory-rich, adaptive spaces?
    \item Finally, could you share a source of inspiration for multi-sensory biophilic design—a book, movie, place, piece of art, or any other experience that has influenced your thinking in this area? Any additional thoughts or anything else that you would like to mention?
\end{enumerate}


